\documentclass[17pt, aspectratio=169]{beamer}

% --- Paquetes de Idioma y Estética ---
\usepackage[utf8]{inputenc}
\usepackage[T1]{fontenc}
\usepackage[spanish]{babel}
\usepackage{xcolor}
\usepackage{tcolorbox}
\tcbuselibrary{skins}

% --- Definición de Colores ICMG ---
\definecolor{ICMGBlue}{HTML}{002366}
\definecolor{ICMGGold}{HTML}{996515}
\definecolor{CitationBg}{HTML}{F4F4F4}

% --- Configuración del Tema Beamer ---
\setbeamercolor{palette primary}{bg=ICMGBlue, fg=white}
\setbeamercolor{palette secondary}{bg=ICMGBlue, fg=white}
\setbeamercolor{structure}{fg=ICMGBlue}
\setbeamercolor{title}{fg=ICMGBlue}
\setbeamerfont{title}{series=\bfseries}

\setbeamertemplate{navigation symbols}{}

\setbeamertemplate{footline}{
  \leavevmode%
  \hbox{%
  \begin{beamercolorbox}[wd=.5\paperwidth,ht=2.25ex,dp=1ex,left]{author in head/foot}%
    \hspace*{2ex} Clase 4: Vida Cristiana
  \end{beamercolorbox}%
  \begin{beamercolorbox}[wd=.5\paperwidth,ht=2.25ex,dp=1ex,right]{date in head/foot}%
    Iglesia Cristiana Manantial de la Gracia \hspace*{2ex} \insertframenumber{} / \inserttotalframenumber \hspace*{2ex} 
  \end{beamercolorbox}}%
  \vspace{2pt}
}

% --- Formato para Citas Bíblicas ---
\newtcolorbox{citaBiblicaSlide}[1]{
  enhanced,
  colback=CitationBg,
  colframe=ICMGGold,
  leftrule=3pt,
  rightrule=0pt,
  toprule=0pt,
  bottomrule=0pt,
  arc=0mm,
  fonttitle=\bfseries\color{ICMGBlue},
  title={#1},
  boxrule=0pt,
  italic
}

\title{Clase 4: Vida Cristiana}
%\subtitle{Clase 1: El Fundamento de Nuestra Identidad}
\author{Iglesia Cristiana Manantial de la Gracia}
\date{Serie: El Bautismo}

% --- Configuración de Diapositiva de Título Intermedia ---
\AtBeginSection[]{
  \begin{frame}[plain] % 'plain' quita el pie de página para resaltar el título
    \centering
    \vfill
    {\color{ICMGGold}\rule{0.6\textwidth}{1.5pt}} % Línea decorativa superior
    \vspace{0.4cm}
    
    {\color{ICMGBlue}\scshape\large SECCIÓN \thesection \par} % "SECCIÓN 1", etc.
    \vspace{0.2cm}
    {\color{ICMGBlue}\Huge\bfseries \insertsectionhead \par} % Nombre de la sección
    
    \vspace{0.4cm}
    {\color{ICMGGold}\rule{0.6\textwidth}{1.5pt}} % Línea decorativa inferior
    \vfill
  \end{frame}
}

\begin{document}

\begin{frame}[plain]
    \titlepage
    \centering
    \color{ICMGGold}\rule{0.5\textwidth}{1pt}
\end{frame}

\section{El Cuerpo de Cristo \\ 1 Cor. 12:12-26}

\begin{frame}{El Cuerpo de Cristo}
    \begin{itemize}
        \item Adopci\'on a una familia (1 Cor. 12:12-13).
    \end{itemize}
\end{frame}

\begin{frame}{El Cuerpo de Cristo}
    \begin{itemize}
        \item Adopci\'on a una familia (1 Cor. 12:12-13).
        \item Dones y contentamiento (1 Cor. 12:14-20).
    \end{itemize}
\end{frame}

\begin{frame}{El Cuerpo de Cristo}
    \begin{itemize}
        \item Adopci\'on a una familia (1 Cor. 12:12-13).
        \item Dones y contentamiento (1 Cor. 12:14-20).
        \item Amor y uni\'on (1 Cor. 12:21-26).
    \end{itemize}
\end{frame}

\section{El Cuerpo de Cristo \\ Efesios 4:11-16}

\begin{frame}{El Cuerpo de Cristo}
    \begin{itemize}
        \item Dios capacita (Efesios 4:11-12).
    \end{itemize}
\end{frame}

\begin{frame}{El Cuerpo de Cristo}
    \begin{itemize}
        \item Dios capacita (Efesios 4:11-12).
        \item Edificaci\'on mutua (Efesios 4:13).
    \end{itemize}
\end{frame}

\begin{frame}{El Cuerpo de Cristo}
    \begin{itemize}
        \item Dios capacita (Efesios 4:11-12).
        \item Edificaci\'on mutua (Efesios 4:13).
        \item La Gran Comisi\'on (Efesios 4:14-16)
    \end{itemize}
\end{frame}

\section{¿Qué hago?}
\begin{frame}{¿Qué hago?}
    \begin{itemize}
        \item Glorificar a Dios, conoci\'endole.
    \end{itemize}
        \begin{citaBiblicaSlide}{Catecismo de la Nueva Ciudad, Pregunta 6}
        \textbf{¿Cómo podemos glorificar a Dios?}\\
Glorificamos a Dios disfrutándolo, amándolo, confiando en Él y obedeciendo Su voluntad, Sus mandamientos y Su ley.
    \end{citaBiblicaSlide}
\end{frame}

%\begin{frame}{¿Qué es el bautismo?}
    %\begin{itemize}
        %\item Medio de Gracia (no salva).
        %\item Reconocer que la Biblia es la única %autoridad y estándar infalible.
    %\end{itemize}
    
    %\begin{citaBiblicaSlide}{Juan 17:3}
        %"Y esta es la vida eterna: que te conozcan a ti, %el único Dios verdadero, y a Jesucristo, a quien %has enviado."
    %\end{citaBiblicaSlide}
    
%\end{frame}

\begin{frame}{Conclusi\'on}
    \begin{itemize}
        \item No \textbf{tienes} que hacer algo, tu redenci\'on está pagada.
        \item El gozo en el Señor nos llevar\'a a conocerle y obedecerle (arrepentimiento).
    \end{itemize}
\end{frame}

\begin{frame}{Conclusi\'on}
    \begin{citaBiblicaSlide}{Catecismo de la Nueva Ciudad, Pregunta 44}
    \textbf{¿Qué es el bautismo?}\\
El bautismo es el lavamiento con agua en el nombre del Padre, del Hijo y del Espíritu Santo; representa y sella nuestra adopción en Cristo, nuestro lavamiento del pecado y nuestro compromiso a pertenecer al Señor y a Su iglesia.
    \end{citaBiblicaSlide}
\end{frame}

\end{document}